\documentclass[11pt]{beamer}
\usepackage[utf8]{inputenc}
\usepackage[T1,T2A]{fontenc}
\usepackage[russian]{babel}
\usepackage{color}
\usepackage{listings}
\usepackage{calc}
\usepackage{graphicx}
\usepackage{epstopdf}
\usepackage{hyperref}
\hypersetup{unicode,colorlinks}
\usepackage{csquotes}
\usepackage{upquote}
\usepackage{cprotect}
\usetheme[progressbar=head,numbering=fraction,block=fill]{metropolis}
\usepackage{dejavu}
\usepackage{etoolbox}
\usepackage{bm}
\usepackage{minted}
%\usepackage[defaultsans]{droidsans}
%\usepackage[defaultmono]{droidmono}
%\makeatletter
%\input droid/t1fdm.fd
%\makeatother

\makeatletter
\DeclareRobustCommand*{\bfseries}{%
    \not@math@alphabet\bfseries\mathbf
    \fontseries\bfdefault\selectfont
    \boldmath
}
\makeatother

\makeatletter
\chardef\straightquote@code=\catcode`'
\chardef\backquote@code=\catcode``
\catcode`'=\active \catcode``=\active
\patchcmd{\@noligs}
{\textasciigrave}
{\fixedtextasciigrave}
{}{}
\newcommand{\fixedtextasciigrave}{%
    \makebox[.5em]{\fontencoding{TS1}\fontfamily{fvs}\selectfont\textasciigrave}% Vera Sans
}
\catcode\lq\'=\straightquote@code
\catcode\lq\`=\backquote@code
\makeatletter

\lstset{
    language=Haskell,
    inputencoding=utf8,
    extendedchars=true,
    breaklines=true,
    escapeinside=!!,
    tabsize=4,
    breakatwhitespace=true,
    keepspaces=true,
    basicstyle=\ttfamily,
    commentstyle=\color{gray}\itshape,
    keywordstyle=
    % keywordstyle=\color{blue}, % TODO https://tex.stackexchange.com/questions/415777/avoid-highlighting-keywords-following-certain-words-in-listings
    % addkeywords={Type}
    % deletekeywords={Monad, String, return, array, show},
    % otherkeywords={::, ->}
}
\lstset{
    literate={а}{{\selectfont\char224}}1
    {б}{{\selectfont\char225}}1
    {в}{{\selectfont\char226}}1
    {г}{{\selectfont\char227}}1
    {д}{{\selectfont\char228}}1
    {е}{{\selectfont\char229}}1
    {ё}{{\"e}}1
    {ж}{{\selectfont\char230}}1
    {з}{{\selectfont\char231}}1
    {и}{{\selectfont\char232}}1
    {й}{{\selectfont\char233}}1
    {к}{{\selectfont\char234}}1
    {л}{{\selectfont\char235}}1
    {м}{{\selectfont\char236}}1
    {н}{{\selectfont\char237}}1
    {о}{{\selectfont\char238}}1
    {п}{{\selectfont\char239}}1
    {р}{{\selectfont\char240}}1
    {с}{{\selectfont\char241}}1
    {т}{{\selectfont\char242}}1
    {у}{{\selectfont\char243}}1
    {ф}{{\selectfont\char244}}1
    {х}{{\selectfont\char245}}1
    {ц}{{\selectfont\char246}}1
    {ч}{{\selectfont\char247}}1
    {ш}{{\selectfont\char248}}1
    {щ}{{\selectfont\char249}}1
    {ъ}{{\selectfont\char250}}1
    {ы}{{\selectfont\char251}}1
    {ь}{{\selectfont\char252}}1
    {э}{{\selectfont\char253}}1
    {ю}{{\selectfont\char254}}1
    {я}{{\selectfont\char255}}1
    {А}{{\selectfont\char192}}1
    {Б}{{\selectfont\char193}}1
    {В}{{\selectfont\char194}}1
    {Г}{{\selectfont\char195}}1
    {Д}{{\selectfont\char196}}1
    {Е}{{\selectfont\char197}}1
    {Ё}{{\"E}}1
    {Ж}{{\selectfont\char198}}1
    {З}{{\selectfont\char199}}1
    {И}{{\selectfont\char200}}1
    {Й}{{\selectfont\char201}}1
    {К}{{\selectfont\char202}}1
    {Л}{{\selectfont\char203}}1
    {М}{{\selectfont\char204}}1
    {Н}{{\selectfont\char205}}1
    {О}{{\selectfont\char206}}1
    {П}{{\selectfont\char207}}1
    {Р}{{\selectfont\char208}}1
    {С}{{\selectfont\char209}}1
    {Т}{{\selectfont\char210}}1
    {У}{{\selectfont\char211}}1
    {Ф}{{\selectfont\char212}}1
    {Х}{{\selectfont\char213}}1
    {Ц}{{\selectfont\char214}}1
    {Ч}{{\selectfont\char215}}1
    {Ш}{{\selectfont\char216}}1
    {Щ}{{\selectfont\char217}}1
    {Ъ}{{\selectfont\char218}}1
    {Ы}{{\selectfont\char219}}1
    {Ь}{{\selectfont\char220}}1
    {Э}{{\selectfont\char221}}1
    {Ю}{{\selectfont\char222}}1
    {Я}{{\selectfont\char223}}1
    {`}{{{\fixedtextasciigrave}}}1
}

% \usemintedstyle{solarized-light}
\newminted[haskell]{haskell}{
    escapeinside=``,
    beameroverlays=true,
    texcomments=true,
    autogobble=true,
    fontsize=\small
}
\newmintinline[haskinline]{haskell}{
    escapeinside=``,
    beameroverlays=true,
    texcomments=true
}

\author{Алексей Романов}
\subtitle{Функциональное программирование на Haskell}
%\logo{}
\institute{МИЭТ}
\subject{Функциональное программирование на Haskell}
%\setbeamercovered{transparent}
%\setbeamertemplate{navigation symbols}{}


% Символы рамки в отчёте Hedgehog.
\DeclareUnicodeCharacter{2503}{\makebox[.6em]{\rule[-.2ex]{.6pt}{1.7ex}}}
\DeclareUnicodeCharacter{2502}{\makebox[.6em]{\rule[-.2ex]{.3pt}{1.7ex}}}
\DeclareUnicodeCharacter{2501}{\rule[.5ex]{.6em}{.6pt}}

\title{Лекция 12: тестирование свойств}

\begin{document}
\begin{frame}[plain]
  \maketitle
\end{frame}

\begin{frame}[fragile]
  \frametitle{Привычный подход к тестированию}
  \begin{itemize}
    \item Как обычно мы тестируем написанные функции?
          \pause
    \item Вычисляем функцию на каких-то конкретных аргументах и сравниваем результат с ожидаемым.
    \item Плюсы:
          \pause
          \begin{itemize}
            \item Тесты легко понять.
            \item Можно сделать даже без помощи библиотек.
            \item Но при этом библиотеки для написания таких тестов есть практически для всех языков.
            \item Для чистых функций (в Haskell почти для всех!) кроме аргументов и результата функции проверять и нечего.
          \end{itemize}
          \pause
    \item TDD (Test Driven Development) \pause "--- сначала тесты, потом реализация.
  \end{itemize}
\end{frame}

\begin{frame}[fragile]
  \frametitle{Недостатки этого подхода}
  \begin{itemize}
    \item А какие у этого подхода минусы?
          \pause
    \item Сколько нужно тестов? Всегда можно добавить ещё.
          \pause
          Разве что если функция на \haskinline|Bool|...
          \pause
    \item Нужно искать пограничные случаи (т.е. те, в которых ошибка более вероятна) самому. \\
          Хотя есть часто встречающиеся: \pause \haskinline|0|, пустой список, \haskinline|INT_MIN|...
          \pause
    \item Для каких-то аргументов мы можем не знать правильного результата (и не иметь независимого от тестируемой функции способа его найти).
          \pause
    \item Есть ситуации, когда программу или функцию в ней намеренно пытаются сломать. \pause Например, злобный преподаватель.
  \end{itemize}
\end{frame}

\begin{frame}[fragile]
  \frametitle{Свойства вместо примеров}
  \begin{itemize}
    \item Тестирование свойств (property-based testing) помогает находить ошибки на множестве автоматически сгенерированных примеров.
    \item Основная идея: вместо результата на конкретных аргументах мы описываем свойства, которые должны выполняться \emph{для всех аргументов} подходящего типа. \pause Или не совсем для всех, но об этом потом.
          \pause
    \item Эти свойства проверяем на случайных аргументах (скажем, 100 за раз по умолчанию).
  \end{itemize}
\end{frame}

\begin{frame}[fragile]
  \frametitle{Что означает успешная проверка?}
  \begin{itemize}
    \item В процессе разработки тесты запускаются далеко не один раз, так что каждое свойство будет проверено в сотнях или тысячах случаев.
          \pause
    \item Успешная проверка не доказывает свойство: результат зависит от выбора генераторов и числа тестов.
          \pause
    \item Важные пограничные случаи всё равно нужно учитывать при построении генераторов.
  \end{itemize}
\end{frame}

\begin{frame}[fragile]
  \frametitle{Простейший пример}
  \begin{itemize}
    \item Будем использовать библиотеку \hackage{hedgehog}:
          \begin{haskellsmall}
            import Hedgehog
            import qualified Hedgehog.Gen as Gen
            import qualified Hedgehog.Range as Range

            prop_square :: Property
            prop_square = property $ do
              (x :: Integer) <- forAll $ Gen.integral $
                Range.linear (-1000) 1000
              assert (x * x >= 0)

            ghci> check prop_square
          \end{haskellsmall}
          \begin{ghci}
            True
          \end{ghci}
    \item Проверили, что квадраты \haskinline|Integer| неотрицательны (100 тестов).
  \end{itemize}
\end{frame}

\begin{frame}[fragile]
  \frametitle{А если использовать Int?}
  \begin{itemize}
    \item Верно ли то же свойство для \haskinline|Int|?
          \begin{haskellsmall}
            prop_squareInt :: Property
            prop_squareInt = property $ do
              x <- forAll $ Gen.int Range.constantBounded
              diff (x * x) (>=) 0
            ghci> check prop_squareInt
          \end{haskellsmall}
          \pause
          {\footnotesize Фрагмент отчёта (64-битный \haskinline|Int|):}
          \begin{ghcismall}
                after 2 tests and 27 shrinks.
            6 ┃   x <- forAll $ Gen.int Range.constantBounded
              ┃   │ 3013262616387239886
            7 ┃   diff (x * x) (>=) 0
              ┃   ^^^^^^^^^^^^^^^^^^^
              ┃   │ ━━━ Failed (- lhs) (+ rhs) ━━━
              ┃   │ - -7464790805971924540
              ┃   │ + 0
            False
          \end{ghcismall}
  \end{itemize}
\end{frame}

\begin{frame}[fragile]
  \frametitle{Переполнение и минимизация}
  \begin{itemize}
    \item \href{https://hackage.haskell.org/package/hedgehog/docs/Hedgehog.html#v:diff}{\haskinline|diff a rel b|} проверяет \haskinline|rel a b|. При ошибке показывает оба значения: \haskinline|-| обозначает левое, \haskinline|+| "--- правое.
    \item Контрпример уменьшен, но не до минимума. Для 64-битного \haskinline|Int| есть, например:
          \begin{haskellsmall}
            ghci> let x = 3037000500 :: Int
            ghci> x * x
          \end{haskellsmall}
          \pause
          \begin{ghcismall}
            -9223372036709301616
          \end{ghcismall}
    \item Из-за переполнения свойство \haskinline|x * x >= 0| может менять истинность при уменьшении \haskinline|x|. Shrinking не гарантирует самого маленького контрпримера.
  \end{itemize}
\end{frame}

\begin{frame}[fragile]
  \frametitle{Ещё один пример}
  \begin{itemize}
    \item Первый элемент списка равен последнему в развёрнутом списке?
          \begin{haskellsmall}
            prop_head :: Property
            prop_head = property $ do
              xs <- forAll $ Gen.list (Range.linear 0 100) $
                Gen.int (Range.linear (-100) 100)
              head xs === last (reverse xs)
            ghci> check prop_head
          \end{haskellsmall}
          \pause
          \begin{ghcismall}
            Prelude.head: empty list
            False
          \end{ghcismall}
          {\footnotesize Вывод сокращён; \haskinline|head| даёт предупреждение \ghcinline|x-partial|.}
    \item Как исправить? \pause Отбросим пустые списки перед вызовом \haskinline|head|:
          \begin{haskellsmall}
            if null xs then discard else success
          \end{haskellsmall}
          Или генерировать непустые списки: \haskinline|Range.linear 1 100|.
  \end{itemize}
\end{frame}

\begin{frame}[fragile]
  \frametitle{Основные типы и функции}
  \begin{itemize}
    \item \haskinline|Gen a|: генератор случайных значений типа \haskinline|a|.
    \item \haskinline|Property|: свойство, которое можно проверить.
    \item Тело свойства "--- в монаде \haskinline|PropertyT IO| (типы упрощены):
          \begin{haskellsmall}
            property :: PropertyT IO () -> Property
            forAll :: Show a => Gen a -> PropertyT IO a
            assert :: Bool -> PropertyT IO ()
            (===) :: (Eq a, Show a) =>
                     a -> a -> PropertyT IO ()
            check :: Property -> IO Bool
          \end{haskellsmall}
    \item \haskinline|forAll| генерирует аргумент; \haskinline|assert| проверяет условие; \haskinline|===| сравнивает значения.
    \item \haskinline|discard| пропускает пример. Слишком много пропусков "--- проблема с генератором.
  \end{itemize}
\end{frame}

\begin{frame}[fragile]
  \frametitle{Генераторы значений}
  \begin{itemize}
    \item Примеры базовых генераторов:
          \begin{haskell}
            Gen.int :: Range Int -> Gen Int
            Gen.element :: [a] -> Gen a
          \end{haskell}
    \item Из простых генераторов можно строить более сложные, пользуясь тем, что \haskinline|Gen| "--- монада. Например:
          \begin{haskell}
            genPair :: Gen a -> Gen b -> Gen (a, b)
            genPair ga gb =!\pause! liftA2 (,) ga gb

            genMaybe :: Gen a -> Gen (Maybe a)
            genMaybe ga = do is_just <- Gen.bool
                            if is_just 
                              then!\pause! fmap Just ga
                              else!\pause! pure Nothing
          \end{haskell}
  \end{itemize}
\end{frame}

\begin{frame}[fragile]
  \frametitle{Композиция и уменьшение}
  \begin{itemize}
    \item Вариант через \haskinline|do| генерирует пары так же:
          \begin{haskellsmall}
            genPairM :: Gen a -> Gen b -> Gen (a, b)
            genPairM ga gb = do
              x <- ga
              y <- gb
              pure (x, y)
          \end{haskellsmall}
          \pause
    \item Начав уменьшать \haskinline|y| в контрпримере, Hedgehog уже не возвращается к \haskinline|x|. С \haskinline|liftA2| можно чередовать уменьшение компонент.
    \item То же ограничение "--- у нескольких \haskinline|forAll| в \haskinline|do|.
    \item В \hackage{falsify} этого ограничения нет: он уменьшает внутренние случайные решения.
  \end{itemize}
\end{frame}

\begin{frame}[fragile]
  \frametitle{Генераторы значений}
  \begin{itemize}
    \item Для рекурсивных типов наивный генератор:
          \begin{haskelltiny}
            import Control.Applicative (liftA3)
            data Tree a = Empty | Node a (Tree a) (Tree a)
              deriving (Eq, Show)

            treeOf :: Gen a -> Gen (Tree a)
            treeOf g = Gen.choice [pure Empty,
                         liftA3 Node g (treeOf g) (treeOf g)]
          \end{haskelltiny}
          но это определение плохое: оно даёт очень маленькие деревья (почему?).
    \item Правильнее:
          \begin{haskelltiny}
            treeOf g = Gen.recursive Gen.choice
              [pure Empty]
              [liftA3 Node g (treeOf g) (treeOf g)]
          \end{haskelltiny}
          \haskinline|Gen.recursive| уменьшает размер при рекурсии и исключает её при малом размере.
  \end{itemize}
\end{frame}

\begin{frame}[fragile]
  \frametitle{Почему деревья маленькие?}
  \begin{itemize}
    \item В первом варианте \haskinline|treeOf| каждый узел с вероятностью $1/2$ "--- лист, иначе у него два потомка.
    \item Это ветвящийся процесс (Гальтона-Ватсона) со средним числом потомков $2 \cdot 1/2 = 1$ "--- \emph{критический} случай.
          \pause
    \item Дерево конечно с вероятностью 1, но его \emph{ожидаемый} размер бесконечен: почти всегда оно крошечное, а изредка генерация не завершается за разумное время.
    \item Отсюда и параметр размера в Hedgehog (\haskinline|Gen.sized|, \haskinline|Gen.resize|): размер должен убывать при рекурсии. \haskinline|Gen.recursive| делает это за нас.
  \end{itemize}
\end{frame}

\begin{frame}[fragile]
  \frametitle{Минимизация контрпримеров}
  \begin{itemize}
    \item Первый найденный контрпример может быть большим.
    \item Библиотека ищет более простые аргументы, дающие ту же ошибку.
    \item В Hedgehog уменьшение (shrinking) встроено в генератор: отдельная функция обычно не нужна.
    \item Например, генератор непустых списков уменьшает контрпример, сохраняя непустоту.
    \item Это не гарантирует минимального контрпримера: качество уменьшения зависит от того, как построен генератор.
    \item Подробнее: \href{https://hypothesis.works/articles/integrated-shrinking/}{Integrated vs type based shrinking}.
  \end{itemize}
\end{frame}

\begin{frame}[fragile]
  \frametitle{Воспроизводимость}
  \begin{itemize}
    \item Если мы нашли контрпример (например, с десятого запуска) и изменили реализацию, нужно проверить, что этот случай исправился.
    \item Для этого можно добавить обычный тест с напечатанными аргументами.
    \item Конечно, это удобно делать только тогда, когда минимизация более-менее удалась.
    \item Кроме того, для некоторых типов результат \haskinline|show x| недостаточен для точного воспроизведения \haskinline|x| (придумаете пример?)
    \item В Hedgehog печатается команда с \href{https://hackage.haskell.org/package/hedgehog/docs/Hedgehog.html#v:recheckAt}{\haskinline|recheckAt|} для воспроизведения по seed и пути уменьшения.
  \end{itemize}
\end{frame}

\begin{frame}[fragile]
  \frametitle{Дополнительные уменьшения для своих типов}
  \begin{itemize}
    \item В \haskinline|genPair ga gb| уменьшение обеих компонент уже задано генераторами \haskinline|ga| и \haskinline|gb|.
    \item Можно добавить свои варианты уменьшения с помощью \haskinline|Gen.shrink|:
          \begin{haskellsmall}
            shrinkTree :: Tree a -> [Tree a]
            shrinkTree Empty = []
            shrinkTree (Node _ l r) = [Empty, l, r]

            genTree :: Gen a -> Gen (Tree a)
            genTree g = Gen.shrink shrinkTree (treeOf g)
          \end{haskellsmall}
    \item Дополнительно пробуем пустое дерево или поддерево.
    \item Экземпляр класса не нужен: генератор явно передаём в \haskinline|forAll|.
  \end{itemize}
\end{frame}

\begin{frame}[fragile]
  \frametitle{Часто встречающиеся виды свойств}
  \begin{itemize}
    \item Самое тривиальное свойство "--- функция даёт результат, а не выкидывает исключение.
    \item Сверка с другой реализацией той же функции.
    \item Сериализация и разбор: \haskinline|decode (encode x) == Right x| "--- закодированное значение можно восстановить.
    \item Алгебраические свойства: идемпотентность, коммутативность, ассоциативность и так далее.
    \item Законы классов типов.
  \end{itemize}
\end{frame}

\begin{frame}[fragile]
  \frametitle{Генераторы функций}
  \begin{itemize}
    \item Для тестирования функций высшего порядка нужны случайные функции.
    \item Для Hedgehog это предоставляет пакет \hackage{hedgehog-fn}, модуль \href{https://hackage.haskell.org/package/hedgehog-fn/docs/Hedgehog-Function.html}{\haskinline|Hedgehog.Function|}:
          \begin{haskellsmall}
            import Hedgehog.Function
              (Arg, Vary, Fn, fn, apply)

            fn :: (Arg a, Vary a) => Gen b -> Gen (Fn a b)
            apply :: Fn a b -> a -> b
          \end{haskellsmall}
    \item \haskinline|Fn a b| представляет функцию \haskinline|a -> b| в виде, который можно показать и уменьшить.
    \item Экземпляры \haskinline|Arg a| и \haskinline|Vary a| нужны для генерации функций с аргументами типа \haskinline|a|; для стандартных типов они уже есть.
  \end{itemize}
\end{frame}

\begin{frame}[fragile]
  \frametitle{Свойства функций}
  \begin{itemize}
    \item Проверим ассоциативность композиции для функций \haskinline|Int -> Int|:
          \begin{haskellsmall}
            prop_compose :: Property
            prop_compose = property $ do
              let gi = Gen.int (Range.linear (-100) 100)
              (fg, fh, fk, x) <- forAll $
                (,,,) <$> fn gi <*> fn gi <*> fn gi <*> gi
              let (g, h, k) = (apply fg, apply fh, apply fk)
              ((g . h) . k) x === (g . (h . k)) x
            ghci> check prop_compose
          \end{haskellsmall}
          \pause
          \begin{ghcismall}
            True
          \end{ghcismall}
    \item У функций нет \haskinline|Eq|: сравниваем результаты на аргументе \haskinline|x|. Функция \haskinline|apply| извлекает обычную функцию из \haskinline|Fn|.
  \end{itemize}
\end{frame}

\begin{frame}[fragile]
  \frametitle{Закон класса как свойство}
  \begin{itemize}
    \item Например, проверим ассоциативность операции \haskinline|<>|:
          \begin{haskellsmall}
            prop_associative :: (Eq a, Show a, Semigroup a)
                             => Gen a -> Property
            prop_associative genA = property $ do
              (x, y, z) <- forAll $
                (,,) <$> genA <*> genA <*> genA
              (x <> y) <> z === x <> (y <> z)
          \end{haskellsmall}
    \item Генератор определяет, какой экземпляр класса проверяем. Например, для списков \haskinline|<>| "--- конкатенация:
          \begin{haskellsmall}
            ghci> check $ prop_associative $
                    Gen.list (Range.linear 0 20) $
                      Gen.int (Range.linear (-100) 100)
          \end{haskellsmall}
          \pause
          \begin{ghcismall}
            True
          \end{ghcismall}
  \end{itemize}
\end{frame}

\begin{frame}[fragile]
  \frametitle{Тестирование законов классов}
  \begin{itemize}
    \item Библиотека \hackage{hedgehog-classes} содержит свойства, проверяющие законы для множества классов.
    \item Например, для \haskinline|Eq|:
          \begin{haskellsmall}
            import Hedgehog.Classes

            eqLaws :: (Eq a, Show a) => Gen a -> Laws
            lawsCheck :: Laws -> IO Bool

            ghci> lawsCheck $ eqLaws $
                    Gen.int (Range.linear (-100) 100)
          \end{haskellsmall}
          \pause
    \item Для проверки законов \haskinline|Functor/Applicative/Monad/...| передаём генератор структуры, например:
          \begin{haskellsmall}
            ghci> lawsCheck (functorLaws Gen.maybe)
          \end{haskellsmall}
  \end{itemize}
\end{frame}

\begin{frame}[fragile]
  \frametitle{Свойства с IO}
  \begin{itemize}
    \item Проверим запись и чтение \haskinline|IORef| из \href{https://hackage.haskell.org/package/base/docs/Data-IORef.html}{\haskinline|Data.IORef|}:
          \begin{haskellsmall}
            import Data.IORef (newIORef, writeIORef, readIORef)

            prop_ioRef :: Property
            prop_ioRef = property $ do
              x <- forAll $ Gen.int Range.constantBounded
              y <- evalIO $ do
                ref <- newIORef 0
                writeIORef ref x
                readIORef ref
              y === x
            ghci> check prop_ioRef
          \end{haskellsmall}
          \pause
          \begin{ghcismall}
            True
          \end{ghcismall}
    \item \haskinline|evalIO| выполняет действие \haskinline|IO| внутри свойства. Новый \haskinline|IORef| создаётся для каждого примера, в том числе при уменьшении.
  \end{itemize}
\end{frame}

\begin{frame}[fragile]
  \frametitle{Альтернативы: QuickCheck}
  \begin{itemize}
    \item \hackage{QuickCheck}: класс \haskinline|Arbitrary| задаёт генератор \haskinline|arbitrary| и отдельную функцию уменьшения \haskinline|shrink|.
    \item Для стандартных типов генератор уже есть; можно сразу записать свойство как функцию:
          \begin{haskellsmall}
            import Test.QuickCheck

            prop_reverse :: [Int] -> Property
            prop_reverse xs = reverse (reverse xs) === xs

            ghci> quickCheck prop_reverse
          \end{haskellsmall}
          \begin{ghcismall}
            +++ OK, passed 100 tests.
          \end{ghcismall}
    \item Генераторы можно передавать и явно, но уменьшение задаётся отдельно.
  \end{itemize}
\end{frame}

\begin{frame}[fragile]
  \frametitle{QuickCheck: переполнение Int}
  \begin{itemize}
    \item Проверим неотрицательность квадрата:
          \begin{haskellsmall}
            prop_squareInt :: Int -> Bool
            prop_squareInt x = x * x >= 0
            ghci> quickCheck prop_squareInt
          \end{haskellsmall}
          \pause
          \begin{ghcismall}
            +++ OK, passed 100 tests.
          \end{ghcismall}
    \item Стандартный генератор \haskinline|Int| начинает с малых значений. Возьмём весь диапазон и явно зададим уменьшение:
          \begin{haskellsmall}
            ghci> quickCheck $
                    forAllShrink arbitraryBoundedIntegral shrink
                      prop_squareInt
          \end{haskellsmall}
          \pause
          {\footnotesize Пример вывода (64-битный \haskinline|Int|):}
          \begin{ghcismall}
            *** Failed! Falsified (after 1 test and 64 shrinks):
            3037000500
          \end{ghcismall}
  \end{itemize}
\end{frame}

\begin{frame}[fragile]
  \frametitle{Альтернативы: Falsify}
  \begin{itemize}
    \item \hackage{falsify}: уменьшает внутренние случайные решения; поддерживает уменьшение через зависимые генераторы.
          \begin{haskellsmall*}{texcomments=false}
            {-# LANGUAGE OverloadedStrings #-}
            import Test.Falsify
            import Test.Falsify.Interactive (falsify)
            import qualified Test.Falsify.Generator as Gen
            import qualified Test.Falsify.Predicate as P
            import qualified Test.Falsify.Range as Range
          \end{haskellsmall*}
    \item Свойство имеет тип \haskinline|Property ()|; \haskinline|gen| играет роль \haskinline|forAll| из Hedgehog.
    \item \haskinline|assert| принимает специальный предикат, например \haskinline|P.eq|. Оператор \haskinline|.$| передаёт ему аргумент вместе с подписью для вывода ошибки.
  \end{itemize}
\end{frame}

\begin{frame}[fragile]
  \frametitle{Falsify: пример}
  \begin{itemize}
    \item Та же проверка двойного разворота списка:
          \begin{haskellsmall}
            prop_reverse :: Property ()
            prop_reverse = do
              xs <- gen $ Gen.list (Range.inclusive (0, 100)) $
                Gen.int (Range.inclusive (-100, 100))
              assert $ P.eq .$ ("actual", reverse (reverse xs))
                            .$ ("expected", xs)

            ghci> falsify prop_reverse
          \end{haskellsmall}
          \begin{ghcismall}
            Nothing
          \end{ghcismall}
    \item \haskinline|Nothing| означает, что контрпример не найден; при нарушении свойства получим \haskinline|Just| с описанием ошибки.
  \end{itemize}
\end{frame}

\begin{frame}[fragile]
  \frametitle{Falsify: переполнение Int}
  \begin{itemize}
    \item \haskinline|Range.uniform| охватывает весь диапазон типа и задаёт уменьшение к нулю:
          \begin{haskellsmall}
            prop_squareInt :: Property ()
            prop_squareInt = do
              x <- gen $ Gen.int Range.uniform
              assert $ P.ge .$ ("x * x", x * x)
                            .$ ("zero", 0)
            ghci> falsify prop_squareInt
          \end{haskellsmall}
          \pause
    \item Получим \haskinline|Just| с описанием ошибки. Пример сообщения для 64-битного \haskinline|Int|:
          \begin{ghcismall}
            x * x < zero
            x * x: -9223372036709301616
            zero : 0
          \end{ghcismall}
    \item \haskinline|P.ge| проверяет \haskinline|>=|; подписи помогают понять, что именно оказалось неверным.
  \end{itemize}
\end{frame}

\begin{frame}[fragile]
  \frametitle{Альтернативы: TinyCheck}
  \begin{itemize}
    \item \href{https://github.com/turion/tinycheck}{TinyCheck}: перебирает значения от простых к сложным вместо случайной генерации. Проверки детерминированы; отдельное уменьшение не требуется.
    \item Пример с библиотекой запуска тестов \hackage{tasty}:
          \begin{haskellsmall}
            import Test.Tasty (defaultMain)
            import Test.Tasty.TinyCheck (testProperty)

            main :: IO ()
            main = defaultMain $
              testProperty "reverse . reverse == id" $
                \(xs :: [Int]) -> reverse (reverse xs) == xs
          \end{haskellsmall}
    \item Использует свой класс \haskinline|Arbitrary| для перечисления примеров; экземпляры для стандартных типов уже есть.
  \end{itemize}
\end{frame}

\begin{frame}[fragile]
  \frametitle{TinyCheck: переполнение Int}
  \begin{itemize}
    \item Та же проверка со стандартным перечислением:
          \begin{haskellsmall}
            main :: IO ()
            main = defaultMain $
              testProperty "square" $
                \(x :: Int) -> x * x >= 0
          \end{haskellsmall}
          \pause
          \begin{ghcismall}
            square: OK
              OK, checked 10000 cases
          \end{ghcismall}
    \item Перебираются \haskinline|0, -1, 1, -2, 2, ...|: до значений, вызывающих переполнение, ещё очень далеко.
    \item Добавим большие значения в перечисление явно.
  \end{itemize}
\end{frame}

\begin{frame}[fragile]
  \frametitle{TinyCheck: добавляем большие значения}
  \begin{itemize}
    \item Значения вида $2^k-1$ чередуем со стандартными:
          \begin{haskellsmall}
            import Data.Bits (finiteBitSize)
            import Data.TestCases (TestCases(..), arbitrary)
            import Test.Tasty.TinyCheck (forAllShow, assert)

            ints :: TestCases Int
            ints = TestCases [2^k - 1 | k <- [0 .. bits - 1]]
              <> arbitrary
              where bits = finiteBitSize (0 :: Int)
            main = defaultMain $ testProperty "square" $ do
              x <- forAllShow "x" ints
              assert (x * x >= 0)
          \end{haskellsmall}
          \pause
          {\footnotesize Сокращённый вывод для 64-битного \haskinline|Int|:}
          \begin{ghcismall}
            square: FAIL
              x: 4294967295
              falsified
          \end{ghcismall}
  \end{itemize}
\end{frame}

\begin{frame}[fragile]
  \frametitle{Дополнительное чтение}
  \footnotesize
  \begin{itemize}
    \item \href{https://hypothesis.works/articles/quickcheck-in-every-language/}{QuickCheck in Every Language} Список библиотек тестирования свойств на разных языках на апрель 2016 (от автора \href{https://hypothesis.works/}{Hypothesis} для Python).
    \item \href{https://begriffs.com/posts/2017-01-14-design-use-quickcheck.html}{The Design and Use of QuickCheck}
    \item \href{https://teh.id.au/posts/2017/04/23/property-testing-with-hedgehog/}{Property testing with Hedgehog}
    \item \href{http://www.well-typed.com/blog/2019/05/integrated-shrinking/}{Integrated versus Manual Shrinking} Реализация и сравнение мини-версий QuickCheck и Hedgehog.
    \item \href{https://well-typed.com/blog/2023/04/falsify/}{falsify: Hypothesis-inspired shrinking for Haskell} (2023).
    \item \href{https://habr.com/ru/post/434008/}{Как перестать беспокоиться и начать писать тесты на основе свойств}
    \item \href{https://www.hillelwayne.com/post/contract-examples/}{Finding Property Tests}
    \item \href{https://wickstrom.tech/programming/2019/03/02/property-based-testing-in-a-screencast-editor-introduction.html}{Property-Based Testing in a Screencast Editor}
    \item \href{https://austinrochford.com/posts/2014-05-27-quickcheck-laws.html}{Verifying Typeclass Laws in Haskell with QuickCheck}
    \item \href{https://blog.jrheard.com/hypothesis-and-pexpect}{Using Hypothesis and Pexpect to Test High School Programming Assignments}
  \end{itemize}
\end{frame}

\end{document}
